\uuid{QrlX}
\titre{Domaine de définition}
\auteur{bodin ; hanani ; mzari}
\organisation{exo7}
\licence{CC BY-NC-SA 4.0}
\datecreate{2020-01-20}
\niveau{L1}
\module{Analyse}
\chapitre{Continuité, limite et étude de fonctions réelles}
\sousChapitre{Autre}
\difficulte{2}
\theme{Notion de fonctions, domaine de définition}
\exo7qcmid{lille-1.continuite.007}

\contenu{
\begin{question}
\qclassification{123.00}
\qtags{motcle=domaine de définition}
Soit $f(x) = \ln(x-1)$ et $g(x) = \sqrt{x+1}$.
Quelles sont les assertions vraies concernant les domaines de définition ? (Rappel : le domaine de définition de $f$ est le plus grand ensemble $D_f \subset \Rr$ sur lequel $f$ est définie.)
\begin{answers}
    \good{$D_f \cup D_g = [-1,+\infty[$.}

    \bad{Pour la composition $f \circ g$, $D_{f\circ g} = [-1,+\infty[$.}

    \bad{Pour la composition $g \circ f$, $D_{g\circ f} = ]1,+\infty[$.}

    \good{Pour la fonction $f \times g$, $D_{f\times g} = ]1,+\infty[$.}
\end{answers}
\begin{explanations}
$D_f = ]1,+\infty[$ ; $D_g = [-1,+\infty[$ ;
$D_f \cup D_g = [-1,+\infty[$ ;
$D_{f \times g} = D_f \cap D_g = ]1,+\infty[$ ;
$D_{f \circ g} = ]1,+\infty[$.
$g \circ f$ n'est pas définie pour $x$ proche de $1$, en fait
$D_{g\circ f} = [1+\frac1e,+\infty[$.
\end{explanations}
\end{question}
}
